\documentclass[11pt]{article}
\usepackage[margin=1in]{geometry}
\usepackage{amsmath,amssymb,amsthm,booktabs}
\usepackage[hidelinks]{hyperref}
\usepackage{xurl}
\setlength{\emergencystretch}{2em}
\newtheorem{theorem}{Theorem}
\newtheorem{lemma}[theorem]{Lemma}
\newtheorem{corollary}[theorem]{Corollary}
\newtheorem{definition}[theorem]{Definition}
\newcommand{\mr}{\operatorname{mr}}
\title{A seventeen-crossing knot with a sharp simple-group bound}
\author{Written proof, finite certificates and verified intervals}
\date{17 November 2026}
\begin{document}\maketitle

An explicit closed four-braid has meridional rank and bridge number four.
A double-transposition epimorphism onto $A_8$ supplies the lower bound.
Its Alexander polynomial and Yamada's theorem show that its genus,
braid index and crossing number are $7,4,17$, respectively. Verified
intervals prove that its knot exterior is hyperbolic. A different
interval certificate gives a nonzero-volume representation of the
actual double branched cover, proving that this cover is not a graph
manifold. The latter argument permits negative tetrahedron shapes
and does not assert hyperbolicity of the closed cover.

The proofs of these statements are separate. In particular, neither
the geometric claims nor a comparison with other equality methods is
a premise of the finite meridional bound. The knot is defined below
by a literal braid and PD, rather than assigned a new census name.
Worldwide priority and exclusion from every previously proved family
are not asserted.

\begin{definition}\label{def:knot}
Let $K_*$ be the closure of the four-braid
\[
 \beta=\sigma_1\sigma_2\sigma_3^{-2}\sigma_2\sigma_3^{-2}
        \sigma_2\sigma_1^2\sigma_2^3\sigma_3^{-1}\sigma_2^3.
\]
The generator convention is the Spherogram braid-closure convention,
bound to the literal PD in Section~\ref{sec:literal}. Reversing all
crossing conventions gives a mirror and preserves the conclusions.
\end{definition}

\section{The finite meridional certificate}

Use $A_8$ on $\{0,\ldots,7\}$ and the quartet
\[
\begin{aligned}
 x_0&=(0\,1)(2\,3),&x_1&=(1\,4)(2\,3),\\
 x_2&=(0\,5)(1\,6),&x_3&=(1\,2)(6\,7).
\end{aligned}
\]
All double transpositions belong to one $A_8$ conjugacy class, of size
$\binom84\cdot3=210$. An odd $S_8$ conjugator can be made even by
multiplying by one transposition factor centralizing its source.

\begin{theorem}\label{thm:sharp}
There is a meridian-class epimorphism
$\pi_1(S^3\setminus K_*)\twoheadrightarrow A_8$ taking meridians to
double transpositions. Moreover,
\[
                         \mr(K_*)=b(K_*)=4.
\]
\end{theorem}
\begin{proof}
The exact literal certificate in Section~\ref{sec:literal} has seventeen
arc images, all double transpositions, and satisfies every relation
$p_op_ap_o=p_b$ at an over/under crossing. The involutions make these
relations independent of the signed Wirtinger convention. Its four
seed images are $x_0,x_1,x_2,x_3$. Exact multiplication enumerates
$20160$ elements in their generated subgroup. Since every image is
even, the image is all of $A_8$.

The lower bound needs only transitivity on the eight points.
For any family of $r$ double transpositions, draw their two
transposition edges per generator. The components of this graph are
exactly the permutation orbits. A transitive family needs at least
seven edges; hence $7\le2r$ and $r\ge4$. Every meridional generating
family maps into the double-transposition class, since meridians are
conjugate up to inversion and these images are involutions.
Thus $\mr(K_*)\ge4$.

The retained four-seed trace colors all seventeen arcs using only the
published Wirtinger coloring move. The bridge--Wirtinger equality
\cite[Definition~1.2 and Theorem~1.3]{bkvv} gives $b(K_*)\le4$.
The elementary bridge presentation gives $\mr(K_*)\le b(K_*)$,
proving equality. Meridional generation alone is not used to infer
the bridge upper bound.
\end{proof}

The target is the ordinary-rank-two group $A_8$, whose specified
meridian class has generating rank four. These are properties of the
finite target, not assertions about ordinary generating rank of the
knot group. They are proved separately in
\path{A8_SIMPLE_QUOTIENTS.tex}, Theorem~\texttt{thm:group}; the
transitivity proof above supplies the necessary lower bound directly.

\section{Genus and minimal crossing number}

\begin{theorem}\label{thm:crossing}
Up to a Laurent unit, the Alexander polynomial of $K_*$ is
\[
\begin{aligned}
\Delta(t)={}&t^{14}-4t^{13}+8t^{12}-8t^{11}+18t^9-38t^8+47t^7\\
           &{}-38t^6+18t^5-8t^3+8t^2-4t+1.
\end{aligned}
\]
In particular $|\Delta(-1)|=117$. The genus is seven, the braid
index is four, and the crossing number is exactly seventeen.
\end{theorem}
\begin{proof}
The checker reconstructs coherent orientations from the literal PD
and forms its Alexander--Fox matrix. At a crossing
$g_b=g_o g_a g_o^{-1}$ the row has coefficients $1-t,t,-1$ on
$o,a,b$; for the reversed conjugation these become $1-t,-1,t$.
A $16\times16$ cofactor presents the Alexander module, after removing
the redundant Wirtinger relation and the constant direction in the
abelianized boundary map. Its determinant is $\Delta$ up to a Laurent
unit. This is also seen directly from the infinite-cyclic-cover chain
complex: the kernel of the generator boundary map has basis
$e_i-e_{17}$, and the cofactor is the relation matrix in that basis.

The determinant has degree at most sixteen. The standard-library
checker evaluates the integer determinant at $t=0,\ldots,16$ by
fraction-free elimination and interpolates the polynomial exactly.
Removing its Laurent unit gives the displayed polynomial. An
independent Burau calculation is a supporting control, not needed by
the canonical replay.

For a genus-$g$ Seifert surface, cutting the infinite cyclic cover
along its lifts gives the classical $2g\times2g$ Seifert-matrix
presentation $tV-V^T$. Thus the Alexander breadth is at most $2g$,
and breadth fourteen forces $g(K_*)\ge7$.
The Seifert surface of this seventeen-crossing closed four-braid has
four Seifert circles and genus $(17-4+1)/2=7$. Therefore $g(K_*)=7$.

Theorem~\ref{thm:sharp} and the four-braid presentation imply that
its braid index is four: bridge number never exceeds braid index.
Yamada's theorem \cite[Theorem~3, printed p.~347]{yamada} states that
the minimal number of Seifert circles equals braid index.
Consequently every diagram $D$ of $K_*$ has $s(D)\ge4$.
Its canonical Seifert genus $g_D$ is at least seven, so
\[
           c(D)=2g_D+s(D)-1\ge14+4-1=17.
\]
The displayed braid attains seventeen crossings.
\end{proof}

\section{The actual manifold certificates}

\begin{theorem}\label{thm:hyperbolic}
The exterior of $K_*$ admits a complete finite-volume hyperbolic
metric. In particular, $K_*$ is a prime hyperbolic knot.
\end{theorem}
\begin{proof}
The topology replay constructs the exterior directly from the literal
PD and checks its saved triangulation including peripheral curves.
The Sage replay certifies a solution in the retained boxes for its
$26$ shapes. It checks that all imaginary parts are positive, that
all edge logarithmic equations have the correct $2\pi i$ lift, and
that both complete-cusp holonomies have lift zero. These are the
usual sufficient gluing conditions for a complete hyperbolic metric.

This step imports the pinned SnapPy~3.3.2/SageMath~10.9 verified
interval algorithms. Shape-solution existence is established with
complex Krawczyk intervals; the remaining inequalities use rigorous
interval arithmetic. The actual replay, version information,
source hashes and rational interval endpoints are retained. A default
standard-library run alone does not perform these analytic checks.
\end{proof}

Let $E$ be that literal exterior, with its actual meridian $m$ and
longitude $l$. Put $Y=\Sigma_2(K_*)$.
The unique nontrivial meridional character sends $m$ to $1$ and $l$
to $0$ in $\mathbb Z/2$. On the boundary its kernel is the lattice
\[
                         \langle 2m,l\rangle.
\]
The primitive lift of $m^2$ is therefore the slope to fill in the
unbranched double cover. This identifies the resulting closed manifold
with the double branched cover.

The replay implements this construction by first declaring the base
orbifold filling $(2,0)$ and then taking its unique cyclic two-sheet
cover. It never guesses a filling on an arbitrary unfilled-cover basis.
The simplified orbifold presentation has four generators $a,b,c,d$,
meridian $d$, a relator $d^2$, and mod-two relator rank three.
The exact character is $a,b,c,d\mapsto1$; the recorded longitude maps
to zero. A second cover construction from this explicit permutation
character gives exactly the same decorated triangulation signature
as the enumerated cover.

SnapPy's retained covering source explicitly lifts the base filling
curve, divides by the primitive intersection-number gcd, adjusts the
orbifold multiplicity, and installs this curve as the new meridian.
The reported coefficient $(1,0)$ is this tracked primitive lift.
In a boundary basis $\widetilde{m^2},\widetilde l$ the covering matrix
is $\operatorname{diag}(2,1)$. A possible longitude shear does not
change the filled slope. The saved filling produces a valid closed
orientable $25$-tetrahedron triangulation with $H_1(Y)=\mathbb Z/117$;
these are additional reconstruction controls.

\begin{theorem}\label{thm:nongraph}
The actual double branched cover $Y=\Sigma_2(K_*)$ is not a graph
manifold.
\end{theorem}
\begin{proof}
Let $M$ be the unfilled one-cusped triangulation underlying the
tracked cover. Its ideal triangulation has $52$ shapes $z_j$, and
write
\[
                    z'_j=(1-z_j)^{-1},\qquad z''_j=(z_j-1)/z_j.
\]
The retained $120$-bit Krawczyk certificate supplies a nondegenerate
solution of the rectangular equations: all shape imaginary intervals
exclude zero, hence $z_j\notin\{0,1\}$. Fourteen shapes have negative
imaginary part. They are permitted in the argument below.

There is an exact algebraic check of the \emph{full} equation system.
A retained unimodular $52\times52$ integer transformation takes the
$52$ edge equations, including their sign constants, to one identity
and $51$ independent equations. Together with the filling equation,
these are exactly the $52$ equations certified by the Krawczyk engine.
The standard-library replay verifies the determinant $\pm1$, all
transformed exponents and constants, and the discarded identity.
Thus the selected certified root satisfies every rectangular equation;
merely finding intervals close to zero would not suffice.

The filling equation's rectangular sign requires care. In a row
$(a_j,b_j,c_j)$ of logarithmic exponents,
\[
 \prod_j z_j^{a_j}(z'_j)^{b_j}(z''_j)^{c_j}
   =(-1)^{\sum c_j}\prod_j z_j^{a_j-c_j}(1-z_j)^{c_j-b_j}.
\]
For the filled slope both the substitution sign and the raw rectangular
right side are $-1$. Their product is $+1$: the actual peripheral
multiplier is exactly one. Its logarithmic lift is $-4\pi i$.
The commuting longitude has a certified logarithm with real part
strictly between $17.47$ and $17.48$, so its multiplier has nonunit
modulus. It is loxodromic and its centralizer is diagonal in the basis
of its two fixed points. A commuting peripheral element of multiplier
one is therefore the identity in $\mathrm{PSL}_2(\mathbb C)$.
The holonomy representation kills the primitive filled slope and
factors through $\pi_1(Y)$.

To identify its volume, Francaviglia's straight pseudo-developing-map
construction associates holonomy to a nondegenerate shape solution and
identifies its volume with the signed tetrahedron sum
\cite[Remark~4.4, p.~433]{francaviglia}. The construction includes
nonparabolic peripheral images
\cite[Corollary~4.14 and Remark~4.15, p.~437]{francaviglia}.
His Dehn-filling compatibility result
\cite[Proposition~4.17, p.~437]{francaviglia} says that the real
representation volume is unchanged when the representation factors
through a filling. Thus the resulting representation
$\rho:\pi_1(Y)\to\mathrm{PSL}_2(\mathbb C)$ has
\[
       \operatorname{Vol}(\rho)=\sum_{j=1}^{52}D(z_j),
\]
where $D$ is the Bloch--Wigner dilogarithm. The Sage/Arb interval
replay certifies
\[
                     10.37<\operatorname{Vol}(\rho)<10.39.
\]
Finally every $\mathrm{PSL}_2(\mathbb C)$ representation of a closed
orientable graph manifold has volume zero
\cite[Lemma~4.10, p.~11 of the retained version]{lty}.
This contradicts the certified positive interval.
No positive-shape triangulation or hyperbolic structure of $Y$ is used.
\end{proof}

\section{Connected sums and bounded comparisons}

\begin{corollary}\label{cor:trefoils}
For every integer $s\ge0$, put $K_s=K_*\#(\#^s3_1)$. Then
\[
                         \mr(K_s)=b(K_s)=4+s.
\]
\end{corollary}
\begin{proof}
Retain the transitive eight-sheet $A_8$ action of the seed. Its meridian
has two moving pairs and four fixed points, so $n=8,c=6,\delta=2$.
Let every meridian of each added trefoil act by the same involution.
Its disconnected branched cover consists of two double covers and
four single-sheet $S^3$ components. The trefoil orbifold presentation
\[
             \langle a,b\mid a^2=b^2=(ab)^3=1\rangle\cong S_3
\]
has even kernel $C_3$. Each double component therefore has
$H_1(-;\mathbb F_3)$ of dimension one.

The connected-sum sphere lifts to six spheres, one per meridian cycle.
They join the existing connected cover to the six trefoil components
in a star tree. Mayer--Vietoris adds two to first homology over
$\mathbb F_3$, with no graph-cycle contribution. After $s$ additions,
the connected cover has $h_s=h_0+2s$, where only $h_0\ge0$ is needed.
The arbitrary-cycle cover inequality
\path{GENERAL_BRANCHED_COVER_RANK_BOUNDS.tex},
Theorem~\texttt{thm:general}, gives
\[
 \mr(K_s)\ge\left\lceil\frac{8-1+h_0+2s}{2}\right\rceil\ge4+s.
\]
Joining a four-bridge position of $K_*$ to $s$ two-bridge trefoil
positions gives the upper bound $b(K_s)\le4+s$.
No seed-cover homology calculation, exact normalized supremum or
unrestricted meridional-rank additivity is assumed.
\end{proof}

\begin{corollary}\label{cor:comparison}
The separate universal reflection obstruction excludes every
rank-at-least-four Coxeter reflection-class quotient of $K_*$,
including infinite targets and singular canonical forms. For every $s\ge0$,
any such quotient of $K_s$ has meridian-class generating rank at most
$3+s$, so a single Coxeter reflection-class quotient cannot supply
the sharp $4+s$ lower bound by its class rank. Theorem~\ref{thm:nongraph}
also excludes the graph-double-cover hypothesis of
the corresponding known equality family.
\end{corollary}
\begin{proof}
The Coxeter statement is the independent companion
\path{UNIVERSAL_17_REFLECTION_OBSTRUCTION.tex},
Corollaries~\texttt{cor:coxeter} and \texttt{cor:sums}. These are
characteristic-zero statements.
The second assertion is exactly the non-graph conclusion above;
the bounded literature comparison is recorded in
\path{NEW_SEVENTEEN_LITERATURE_AUDIT.json}.
Neither comparison is a premise of Theorem~\ref{thm:sharp}.
\end{proof}

\section{Literal finite data}\label{sec:literal}

The PD, in braid-crossing order, is
\[
\begin{gathered}
(4,2,5,1),\ (5,17,6,16),\ (27,6,28,7)\\
(7,28,8,29),\ (17,9,18,8),\ (29,18,30,19)\\
(19,30,20,31),\ (9,21,10,20),\ (2,22,3,21)\\
(22,4,23,3),\ (23,11,24,10),\ (11,25,12,24)\\
(25,13,26,12),\ (31,26,32,27),\ (13,33,14,32)\\
(33,15,0,14),\ (15,1,16,0)
\end{gathered}
\]
Join zero-based PD slots $1$ and $3$ (the second and fourth entries)
through each overpass and sort the
connected edge-label components lexicographically to obtain arcs
$0,\ldots,16$. Their images are
\[
\begin{array}{c|c@{\qquad}c|c}
\text{arc}&\text{permutation}&\text{arc}&\text{permutation}\\\hline
0 & (1\,4)(2\,3) & 9 & (0\,4)(2\,6)\\
1 & (0\,1)(2\,3) & 10 & (0\,4)(6\,7)\\
2 & (0\,4)(2\,3) & 11 & (1\,4)(6\,7)\\
3 & (2\,3)(4\,5) & 12 & (2\,4)(6\,7)\\
4 & (1\,2)(4\,5) & 13 & (1\,2)(6\,7)\\
5 & (0\,5)(1\,2) & 14 & (1\,3)(6\,7)\\
6 & (0\,5)(1\,4) & 15 & (1\,3)(2\,7)\\
7 & (0\,5)(4\,6) & 16 & (1\,6)(2\,3)\\
8 & (0\,5)(1\,6) & &\\
\end{array}
\]
The ordered seed arcs are $(1,0,8,13)$, with images
$(x_0,x_1,x_2,x_3)$. Each triple below means
$(\text{known over},\text{known under},\text{new under})$:
\[
\begin{gathered}
(0,1,2),\ (8,2,3),\ (3,13,14)\\
(14,3,4),\ (4,8,9),\ (9,14,15)\\
(15,9,10),\ (10,4,5),\ (1,10,11)\\
(5,11,12),\ (12,5,6),\ (13,15,16)\\
(16,6,7)
\end{gathered}
\]

\section{Replay and trust boundaries}

The four main files are the present source,
\path{new_seventeen_crossing_knot_inputs.json},
\path{verify_new_seventeen_crossing_knot.py} and
\path{new_seventeen_crossing_knot_verification.json}.
The default checker uses only Python's standard library. It verifies
literal braid wiring, closure, all seventeen permutation relations,
the full finite image, the four-seed trace, the Alexander cofactor,
the mod-two covering character, exact unimodular equation reduction
and the structure of all retained rational interval boxes.
It does not independently establish an analytic root or evaluate
transcendental interval functions.

Run the same checker with \texttt{--topology} in the pinned
SnapPy~3.3.2, Spherogram~2.4.1 and Regina~7.4 environment to rebuild the
actual exterior, orbifold cover, primitive filling and closed
triangulation. Run it separately with \texttt{--sage} in SageMath~10.9
with SnapPy~3.3.2 to recertify the complete cusped structure, the
closed-group shape solution, loxodromic peripheral multiplier and
signed dilogarithm interval. These runs retain their own source hashes;
a replay of an earlier edition is not assurance about a later one.

The trusted imports are distinct: the PD/exterior construction,
the topological covering and peripheral-lifting implementation,
the interval Krawczyk solver, and the Sage/Arb dilogarithm implementation.
Relevant SnapPy source snapshots and exact hashes are retained beside
the primary papers. The cover source explicitly handles orbifold
fillings; it is not covered merely by an unfilled-cell construction
claim. The written reconstruction and volume arguments connect the
software outputs to the stated manifolds.

The certificates are internal exact and verified-interval computations,
with explicit published theorem and software imports. Formal
verification and external specialist review are not claimed.
No closed-cover hyperbolicity, atoroidality, full JSJ decomposition,
unrestricted augmentation-rank gap or global-priority conclusion is
asserted.

\begin{thebibliography}{9}\small
\bibitem{bkvv} R. Blair, A. Kjuchukova, R. Velazquez and P. Villanueva,
Wirtinger systems of generators of knot groups,
Commun. Anal. Geom. 28 (2020), 243--262,
Definition~1.2 and Theorem~1.3, p.~246.
\bibitem{yamada} S. Yamada, The minimal number of Seifert circles equals
the braid index of a link, Invent. Math. 89 (1987), 347--356,
Theorem~3, p.~347. DOI:10.1007/BF01389082.
\bibitem{francaviglia} S. Francaviglia, Hyperbolic volume of
representations of fundamental groups of cusped 3-manifolds,
Int. Math. Res. Not. 2004 (9), 425--459.
DOI:10.1155/S1073792804131619.
\bibitem{lty} F. Luo, S. Tillmann and T. Yang,
Thurston's spinning construction and solutions to the hyperbolic
gluing equations for closed hyperbolic 3-manifolds,
arXiv:1004.2992v1 (2010), Lemma~4.10, p.~11 of the exact retained
author version.
\end{thebibliography}
\end{document}
